\documentclass[11pt]{article}

\usepackage[margin=1in]{geometry}
\usepackage{booktabs}
\usepackage{longtable}
\usepackage{array}
\usepackage{graphicx}
\usepackage{hyperref}

\title{TerrainSense: A YOLO-Formatted Terrain Detection Dataset and SQLite Annotation Database}
\author{Draft prepared from the current workspace dataset}
\date{}

\begin{document}
\maketitle

\begin{abstract}
We present a compact terrain-scene object detection dataset curated in YOLO format and packaged with a normalized SQLite database for reproducible analysis. The dataset contains three splits, four object categories, 2,882 labeled images, and 11,841 object instances suitable for training and evaluation of detection models. In addition to the raw label files, we provide a relational representation that stores image metadata, annotation geometry, class metadata, and split membership. This makes the dataset easier to query, validate, and summarize for downstream experiments.
\end{abstract}

\section{Introduction}
Object detection datasets are often distributed as flat folders of images and text labels. That format is convenient for training, but it is awkward for validation, analysis, and paper writing because aggregate questions require custom scripts. To support reproducible study of the TerrainSense dataset, we convert the YOLO directory structure into a small SQLite database and document the resulting corpus in this paper.

\section{Dataset}
The dataset is organized into train, validation, and test splits with YOLO-format annotations. The current workspace metadata defines four classes: obstacle, person, pothole, and vehical. The dataset root is \texttt{terrain\_dataset/}, with split-specific \texttt{images/} and \texttt{labels/} folders.

\subsection{Dataset Size}
The imported dataset contains 2,303 training images, 288 validation images, and 291 test images. The corresponding label files are present for all images in the imported splits, yielding 11,841 total annotations.

\begin{table}[h]
\centering
\begin{tabular}{lrrr}
\toprule
Split & Images & Label files & Annotations \\
\midrule
Train & 2303 & 2303 & 9528 \\
Valid & 288 & 288 & 1194 \\
Test & 291 & 291 & 1119 \\
\bottomrule
\end{tabular}
\caption{Current dataset counts from the workspace. The annotation totals are generated by the database import script.}
\end{table}

\begin{table}[h]
\centering
\begin{tabular}{lr}
\toprule
Class & Annotation count \\
\midrule
obstacle & 7369 \\
person & 2244 \\
pothole & 748 \\
vehical & 1480 \\
\bottomrule
\end{tabular}
\caption{Class frequency summary computed from the imported YOLO labels.}
\end{table}

\section{Database Design}
The SQLite database stores the dataset in four tables.

\begin{itemize}
\item \textbf{classes}: one row per class id and class name.
\item \textbf{images}: one row per image, including split, filename, image dimensions, size on disk, label path, and label presence.
\item \textbf{annotations}: one row per bounding box in YOLO format, with derived pixel-space coordinates for easier analysis.
\item \textbf{metadata}: small key-value table for dataset-level provenance and summary values.
\end{itemize}

This structure is intentionally simple. SQLite is sufficient here because the goal is a portable, single-file research artifact rather than a multi-user production system. The database can be copied with the dataset and queried directly from Python, the SQLite CLI, or visualization tools.

\section{Methods}
The import pipeline reads each image, finds the matching YOLO text file, parses the five-field object rows, and stores both normalized and pixel-space bounding boxes. Images without labels are retained in the database and explicitly marked. This makes the import faithful to the source dataset while also surfacing data quality issues that would otherwise be hidden.

\subsection{Validation}
The importer rejects malformed label rows that do not contain exactly five fields. It also records images with missing label files, which is useful for auditing and for explaining any mismatch between image counts and annotation counts in the paper.

\section{Results}
The generated \texttt{dataset\_stats.json} file provides split-level counts, class frequencies, and image-size summaries. The current dataset contains 2,882 images and 11,841 annotations, with class frequencies skewed toward obstacle instances.

\section{Discussion}
The main benefit of the database is not storage efficiency; it is reproducibility. A relational representation makes it straightforward to ask questions such as how many annotations each class has, how many images are unlabeled, which splits contain the most crowded scenes, and what the box-size distribution looks like. That supports both paper writing and future dataset maintenance.

\section{Conclusion}
We provide a practical packaging of a YOLO detection dataset into a SQLite-backed research artifact. The approach preserves the original training format while adding queryable structure for analysis, reporting, and quality control.

\end{document}